The 208 lab in the ERB, managed by Steven McDermott, contains a large inventory of equipment and other items that are frequently used by students. However, many students are unfamiliar with the lab and the locations of stored equipment. As a result, students often contact Steven McDermott to determine where specific items are located. These inquiries slow down students as they wait for a response and require the lab manager to repeatedly answer questions that could otherwise be resolved independently. This takes time away from Steven McDermott's other responsibilities and is inefficient for accessing information about the lab's inventory.

The current system also makes it difficult for students to find equipment when they do not know its exact name or terminology. For example, a student may know the functionality of an item without knowing what the item is called, or they may enter a search containing a typographical error. A system that relies on users already knowing the correct terminology limits how effectively students can locate equipment. There is an opportunity to improve this process by providing students with a centralized and accessible way to search the lab's inventory.

Steven McDermott has sponsored our team to develop a website that addresses these issues. The website will allow users to search for equipment and other items in the lab and find information about each item, including its location and appearance. The search functionality will be designed to mitigate common user errors, such as typos and searches based on an item's functionality. Speech-to-text functionality will also allow users to search using their voice, providing a convenient method of accessing the inventory.

In addition to improving access to equipment, the project provides an opportunity to increase awareness of the creative projects developed in Lab 208. Many of these projects are currently unknown to students and others outside of the lab. Providing a space to showcase these projects can help inspire other students, demonstrate the work being done in the lab, and give appropriate recognition to the students who created them.